% GENERATED FILE. Edit canonical lessons, then run npm run generate:books.
% canonical-source-hash: fnv1a64:17ea505c6354722d
% canonical-lessons: GU-C163-nasto, GU-C163-bhojan, GU-C163-jamvu, GU-C163-pirasvu, GU-C163-vahecvu

\chapter{Breakfast and the family meal}
\label{ch:gu-breakfast-and-the-family-meal}
\hlchaptermodality{163}

\begin{chapteropening}
\textbf{By the end of this chapter:} \emph{I can talk about breakfast and the family meal.}

Complete the last lesson of chapter 163: I can talk about breakfast and the family meal.
\end{chapteropening}

\section[\texorpdfstring{\gu{નાસ્તો} (nāsto)}{નાસ્તો (nāsto)}]{\gu{નાસ્તો} (nāsto) — breakfast, snack}

\label{lesson:GU-C163-nasto}



\begin{quote}
Before the new one: say the Gujarati for clay pot, then the Gujarati for corner.
\end{quote}

\subsection*{You'll want to know: \gu{નાસ્તો}}
\textbf{\gu{નાસ્તો}} — \emph{nāsto} — \textquotedblleft{}breakfast, snack\textquotedblright{}.

\begin{quote}
\textbf{\gu{અમે} \gu{સવારે} \gu{નાસ્તો} \gu{કરીએ} \gu{છીએ}.} — \emph{ame savāre nāsto karīe chīe} — \textquotedblleft{}We have breakfast in the morning.\textquotedblright{}
\end{quote}

\begin{cousinweb}
\textbf{\gu{નાસ્તો}} (\emph{nāsto}) is a loan from Persian \emph{nāshta}, \textquotedblleft{}breakfast\textquotedblright{}; EtymDB lists the Persian word. CDIAL, which follows inherited words, has no entry for it.
\end{cousinweb}

\begin{grammarlens}[title={What you've built}]
One more word for everyday life at A2.
\end{grammarlens}

\subsection*{Your turn}
\begin{itemize}
\raggedright
  \item \emph{Say it:} \emph{nāsto}
  \item \emph{Say it:} \emph{nāsto}, once more
  \item \emph{Say it:} \emph{ame savāre nāsto karīe chīe} — We have breakfast in the morning.
  \item \emph{Recall:} say the Gujarati for clay pot, then the Gujarati for corner, then say \emph{nāsto} again
\end{itemize}

\begin{tcolorbox}[breakable,colback=teal!4,colframe=teal!35!black,title={Before you move on}]
What does \gu{નાસ્તો} mean? (\textquotedblleft{}Breakfast, snack\textquotedblright{}.) Say it once more.
\end{tcolorbox}
\section[\texorpdfstring{\gu{ભોજન} (bhojan)}{ભોજન (bhojan)}]{\gu{ભોજન} (bhojan) — meal, food}

\label{lesson:GU-C163-bhojan}



\begin{quote}
Before the new one: say the Gujarati for corner, then the Gujarati for breakfast.
\end{quote}

\subsection*{You'll want to know: \gu{ભોજન}}
\textbf{\gu{ભોજન}} — \emph{bhojan} — \textquotedblleft{}meal, food\textquotedblright{}.

\begin{quote}
\textbf{\gu{આજનું} \gu{ભોજન} \gu{સારું} \gu{છે}.} — \emph{ājnũ bhojan sārũ che} — \textquotedblleft{}Today's meal is good.\textquotedblright{}
\end{quote}

\begin{cousinweb}
\textbf{\gu{ભોજન}} (\emph{bhojan}) is Sanskrit \emph{bhojana}, \textquotedblleft{}eating; food\textquotedblright{}, taken over whole as a learned loan, a tatsama. CDIAL lists only the inherited Old Gujarati \emph{bhoyaṇa}, \textquotedblleft{}food, meal\textquotedblright{}, under that word (CDIAL 9630).
\end{cousinweb}

\begin{grammarlens}[title={What you've built}]
One more word for everyday life at A2.
\end{grammarlens}

\subsection*{Your turn}
\begin{itemize}
\raggedright
  \item \emph{Say it:} \emph{bhojan}
  \item \emph{Say it:} \emph{bhojan}, once more
  \item \emph{Say it:} \emph{ājnũ bhojan sārũ che} — Today's meal is good.
  \item \emph{Recall:} say the Gujarati for corner, then the Gujarati for breakfast, then say \emph{bhojan} again
\end{itemize}

\begin{tcolorbox}[breakable,colback=teal!4,colframe=teal!35!black,title={Before you move on}]
What does \gu{ભોજન} mean? (\textquotedblleft{}Meal, food\textquotedblright{}.) Say it once more.
\end{tcolorbox}
\section[\texorpdfstring{\gu{જમવું} (jamvũ)}{જમવું (jamvũ)}]{\gu{જમવું} (jamvũ) — to have a meal, to dine}

\label{lesson:GU-C163-jamvu}



\begin{quote}
Before the new one: say the Gujarati for breakfast, then the Gujarati for meal.
\end{quote}

\subsection*{You'll want to know: \gu{જમવું}}
\textbf{\gu{જમવું}} — \emph{jamvũ} — \textquotedblleft{}to have a meal, to dine\textquotedblright{}.

\begin{quote}
\textbf{\gu{અમે} \gu{સાંજે} \gu{ઘરે} \gu{જમીએ} \gu{છીએ}.} — \emph{ame sā̃je ghare jamīe chīe} — \textquotedblleft{}We eat at home in the evening.\textquotedblright{}
\end{quote}

\begin{cousinweb}
\textbf{\gu{જમવું}} (\emph{jamvũ}) is a tadbhava, inherited through Prakrit. CDIAL lists it under Sanskrit \emph{jámati}, \textquotedblleft{}eats\textquotedblright{}, and again under \emph{*jimati}, and says it may come from either (CDIAL 5126, 5267).
\end{cousinweb}

\begin{grammarlens}[title={What you've built}]
One more word for everyday life at A2.
\end{grammarlens}

\subsection*{Your turn}
\begin{itemize}
\raggedright
  \item \emph{Say it:} \emph{jamvũ}
  \item \emph{Say it:} \emph{jamvũ}, once more
  \item \emph{Say it:} \emph{ame sā̃je ghare jamīe chīe} — We eat at home in the evening.
  \item \emph{Recall:} say the Gujarati for breakfast, then the Gujarati for meal, then say \emph{jamvũ} again
\end{itemize}

\begin{tcolorbox}[breakable,colback=teal!4,colframe=teal!35!black,title={Before you move on}]
What does \gu{જમવું} mean? (\textquotedblleft{}To have a meal, to dine\textquotedblright{}.) Say it once more.
\end{tcolorbox}
\section[\texorpdfstring{\gu{પીરસવું} (pīrasvũ)}{પીરસવું (pīrasvũ)}]{\gu{પીરસવું} (pīrasvũ) — to serve (food)}

\label{lesson:GU-C163-pirasvu}



\begin{quote}
Before the new one: say the Gujarati for meal, then the Gujarati for to have a meal.
\end{quote}

\subsection*{You'll want to know: \gu{પીરસવું}}
\textbf{\gu{પીરસવું}} — \emph{pīrasvũ} — \textquotedblleft{}to serve (food)\textquotedblright{}.

\begin{quote}
\textbf{\gu{માતા} \gu{થાળીમાં} \gu{ભાત} \gu{પીરસે} \gu{છે}.} — \emph{mātā thāḷīmā̃ bhāt pīrase che} — \textquotedblleft{}Mother serves rice on the plate.\textquotedblright{}
\end{quote}

\begin{cousinweb}
\textbf{\gu{પીરસવું}} (\emph{pīrasvũ}) is a tadbhava. It comes from Sanskrit \emph{pariveṣati}, \textquotedblleft{}attends, serves food\textquotedblright{}, and CDIAL glosses the Gujarati verb \textquotedblleft{}to serve a meal\textquotedblright{} (CDIAL 7888).
\end{cousinweb}

\begin{grammarlens}[title={What you've built}]
One more word for everyday life at A2.
\end{grammarlens}

\subsection*{Your turn}
\begin{itemize}
\raggedright
  \item \emph{Say it:} \emph{pīrasvũ}
  \item \emph{Say it:} \emph{pīrasvũ}, once more
  \item \emph{Say it:} \emph{mātā thāḷīmā̃ bhāt pīrase che} — Mother serves rice on the plate.
  \item \emph{Recall:} say the Gujarati for meal, then the Gujarati for to have a meal, then say \emph{pīrasvũ} again
\end{itemize}

\begin{tcolorbox}[breakable,colback=teal!4,colframe=teal!35!black,title={Before you move on}]
What does \gu{પીરસવું} mean? (\textquotedblleft{}To serve (food)\textquotedblright{}.) Say it once more.
\end{tcolorbox}
\section[\texorpdfstring{\gu{વહેંચવું} (vahẽcvũ)}{વહેંચવું (vahẽcvũ)}]{\gu{વહેંચવું} (vahẽcvũ) — to share out, to divide}

\label{lesson:GU-C163-vahecvu}



\begin{quote}
Before the new one: say the Gujarati for to have a meal, then the Gujarati for to serve.
\end{quote}

\subsection*{You'll want to know: \gu{વહેંચવું}}
\textbf{\gu{વહેંચવું}} — \emph{vahẽcvũ} — \textquotedblleft{}to share out, to divide\textquotedblright{}.

\begin{quote}
\textbf{\gu{બાળકો} \gu{રોટલી} \gu{વહેંચે} \gu{છે}.} — \emph{bāḷko roṭlī vahẽce che} — \textquotedblleft{}The children share the rotli.\textquotedblright{}
\end{quote}

\begin{cousinweb}
\textbf{\gu{વહેંચવું}} (\emph{vahẽcvũ}) is a tadbhava. It comes from Sanskrit \emph{vibhakta}, \textquotedblleft{}divided, distributed\textquotedblright{}, and CDIAL glosses the Gujarati verb \textquotedblleft{}divides, distributes\textquotedblright{} (CDIAL 11808).
\end{cousinweb}

\begin{grammarlens}[title={What you've built}]
One more word for everyday life at A2.
\end{grammarlens}

\subsection*{Your turn}
\begin{itemize}
\raggedright
  \item \emph{Say it:} \emph{vahẽcvũ}
  \item \emph{Say it:} \emph{vahẽcvũ}, once more
  \item \emph{Say it:} \emph{bāḷko roṭlī vahẽce che} — The children share the rotli.
  \item \emph{Recall:} say the Gujarati for to have a meal, then the Gujarati for to serve, then say \emph{vahẽcvũ} again
\end{itemize}

\begin{tcolorbox}[breakable,colback=teal!4,colframe=teal!35!black,title={Before you move on}]
What does \gu{વહેંચવું} mean? (\textquotedblleft{}To share out, to divide\textquotedblright{}.) Say it once more.
\end{tcolorbox}
